\documentclass[sigconf,nonacm]{acmart}

\title{SecureDrop: A Secure Cross-Language Peer-to-Peer File Sharing System}

\author{Monica Stef}
\affiliation{%
  \institution{Queen's University}
}
\author{Ben Leray}
\affiliation{%
  \institution{Queen's University}
}
\author{Anjali Patel}
\affiliation{%
  \institution{Queen's University}
}
\renewcommand{\shortauthors}{Stef, Leray, and Patel}

\begin{document}

\begin{abstract}
SecureDrop is a peer-to-peer file sharing application implemented across
three languages: Go, Python, and Java. The project focuses on secure
discovery, authenticated session establishment, encrypted communication,
verified file transfer, and protected local storage. Each peer maintains
an Ed25519 identity, performs a signed X25519 handshake to derive a fresh
session key, and exchanges application messages over AES-GCM encrypted
channels. The system supports user-approved file transfer, verified
re-sharing using original-owner metadata, and identity key rotation.
Testing in the repository covers protocol helpers, cryptographic
round-trips, metadata persistence, filtered file listing, and selected
security-sensitive update paths. The result is a working interoperable
prototype that demonstrates how secure systems concepts can be implemented
consistently across multiple programming languages.
\end{abstract}

\keywords{peer-to-peer systems, secure file sharing, authenticated key exchange, AES-GCM, Ed25519, X25519, cross-language interoperability}

\maketitle

\section{Introduction}
Secure file sharing is often presented either as a networking problem or
as a cryptography problem. In practice, real systems require both. A
useful file sharing system must allow peers to locate one another,
authenticate remote identities, protect data in transit, verify the
origin of shared content, and preserve security when files are re-shared
between multiple participants.

This project implements these ideas in a cross-language setting. Our
system, SecureDrop, includes Go, Python, and Java clients that speak a
common protocol over TCP. The design goal was not only to transfer files,
but to do so with explicit authentication and integrity guarantees while
keeping the architecture simple enough to reason about and test.

\section{System Overview}
SecureDrop is a localhost/LAN peer-to-peer application in which each
client can listen for incoming connections, connect to known peers,
request file listings, send file requests, and download or re-share
content. The repository contains three client implementations:
\begin{itemize}
\item a Go client,
\item a Python client, and
\item a Java client.
\end{itemize}

Each client maintains a \texttt{shared\_files/} directory for files it is
willing to share and a \texttt{downloads/} directory for received files.
The common command set includes peer listing, remote file listing, file
request, connectivity testing through ping, and identity rotation.

The project specifically implements the following course features:
\begin{itemize}
\item user consent before file transfer,
\item verified file re-sharing,
\item key rotation, and
\item secure local storage for downloaded files.
\end{itemize}

\section{Protocol Design}
\subsection{Handshake}
SecureDrop uses a line-based protocol. During connection setup, each peer
sends a \texttt{HELLO} message containing its user name, Ed25519 public
key, ephemeral X25519 public key, and a signature over the ephemeral
public key. The receiver verifies that signature before proceeding.

After both \texttt{HELLO} messages are exchanged, each side computes a
shared secret using X25519 and derives the session key as the SHA-256 hash
of that shared secret. This provides authenticated session establishment
and fresh per-connection keys. Because the long-term identity key signs
the ephemeral key, the handshake binds the temporary Diffie--Hellman
exchange to a stable identity.

\subsection{Encrypted Messaging}
After the handshake, all application payloads are sent inside encrypted
\texttt{DATA} records. The payload is encrypted with AES-GCM using the
derived session key and a fresh nonce. This gives confidentiality and
integrity for commands such as:
\begin{itemize}
\item \texttt{PING},
\item \texttt{LIST\_REQ} and \texttt{LIST\_RES},
\item \texttt{GET\_REQ} and \texttt{GET\_RES},
\item \texttt{KEY\_UPDATE}, and
\item \texttt{ERROR}.
\end{itemize}

\subsection{File Transfer and Metadata}
When a peer requests a file with \texttt{GET\_REQ}, the sender does not
transfer immediately. Instead, the user is prompted to approve or reject
the request. If approved, the sender returns a \texttt{GET\_RES} payload
containing:
\begin{itemize}
\item the filename,
\item the file contents,
\item the SHA-256 hash of the contents,
\item a signature on that hash,
\item the original owner's public key, and
\item the original owner's name.
\end{itemize}

This metadata allows a receiver to validate both integrity and authorship.
If a downloaded file is later re-shared, the sender can preserve the
original metadata rather than claiming authorship itself.

\section{Security Features}
\subsection{Mutual Authentication}
All three clients authenticate the remote peer during the handshake by
verifying an Ed25519 signature over the remote ephemeral key. In the Go
client, an additional identity consistency check rejects a peer that
reconnects under the same name with a different public key unless a valid
key rotation has already updated trust.

\subsection{Confidentiality and Integrity in Transit}
Application traffic is protected with AES-GCM after session key
derivation. Since AES-GCM is an authenticated encryption scheme, payload
tampering is detected during decryption.

\subsection{Verified Re-Sharing}
Verified re-sharing is one of the most important project features. When a
file originates locally, the sender hashes the contents with SHA-256 and
signs the hash using its Ed25519 private key. When the file originates
from a previous download, the sender reuses the stored origin metadata.
The receiver recomputes the hash, compares it to the transmitted hash, and
verifies the signature using the original owner's public key. This allows
trust in the original source to survive across multiple sharing hops.

\subsection{Secure Local Storage}
Downloaded files are not stored in plaintext. In the Go and Python
clients, files are encrypted before being written to disk and are stored
in the form \texttt{base64(nonce)|base64(ciphertext)}. A separate metadata
file stores the original owner, public key, hash, and signature. This
protects against simple offline disclosure of downloaded content.

The repository shows one implementation difference worth noting: the Go
and Python clients derive a local storage key from the identity name,
while the Java client currently uses a fixed local secret for downloaded
file encryption. The report therefore treats secure local storage as fully
implemented in Go and Python and partially aligned in Java.

\subsection{Key Rotation}
SecureDrop also supports identity rotation. The Go and Python clients send
a \texttt{KEY\_UPDATE} message containing the new public key and a
signature from the old key, allowing peers to verify that the update is
authorized. This preserves trust while allowing an identity to migrate.

The current Java client accepts a simpler \texttt{KEY\_UPDATE} format and
updates the trusted key without verifying an old-key signature. As a
result, signed key rotation is fully implemented in Go and Python and only
partially implemented in Java in the current repository state.

\section{Implementation}
\subsection{Go Client}
The Go client is the most complete reference implementation in the
repository. It supports the full command loop, signed X25519 handshake,
AES-GCM message protection, queued approval for incoming file requests,
verified re-sharing, signed key update handling, and encrypted local
download storage. It also contains repository tests for cryptographic
round-trips, metadata persistence, pending request handling, and verified
key update processing.

\subsection{Python Client}
The Python client mirrors the same protocol structure using the
\texttt{cryptography} library. It supports handshake verification,
encrypted application messages, metadata-aware file transfer, encrypted
download storage, and signed key rotation verification. The Python tests
cover metadata round-trip behavior, encrypted stored-file round trips,
storage format rejection for malformed data, shared file filtering, and
identity-dependent local storage keys.

\subsection{Java Client}
The Java client uses Bouncy Castle primitives for Ed25519 and X25519 and
implements the common TCP message flow, authenticated handshake, AES-GCM
message transport, file approval, and metadata-based verification on
download. The Java implementation interoperates at the handshake and file
transfer level, but the current repository code shows a simplified
approach for key updates and local storage compared to Go and Python.

\section{Testing and Evaluation}
The repository includes both automated and manual testing guidance.
Automated tests validate core helper functions in each language:
\begin{itemize}
\item Go tests cover AES-GCM round-trip encryption, encrypted download
storage round-trip, metadata serialization, hidden metadata files in
shared listings, queued pending requests, and verified key update
handling.
\item Python tests cover metadata save/load, encrypted download
save/load, shared file filtering, invalid stored-file rejection, and
deterministic per-identity storage keys.
\item Java tests cover AES-GCM round-trip encryption, metadata
serialization, hidden metadata files in shared listings, and Ed25519
sign/verify behavior used by handshake helpers.
\end{itemize}

Manual end-to-end tests in the repository exercise peer discovery or
reachability, file listing, user-approved transfer, verified download,
integrity failure on tampering, signature failure on metadata tampering,
key rotation, disconnect handling, and verified re-sharing across peers.

Overall, the test suite gives strong confidence in the protocol helpers
and storage paths, while full three-way interoperability still depends on
manual execution of the clients together.

\section{Limitations}
Although the project successfully demonstrates secure peer-to-peer file
sharing, the current prototype still has limitations.

First, implementation parity is not perfect across all languages. The Go
and Python clients implement stronger versions of key rotation and local
storage protection than the current Java client. Second, trust management
is still minimal; there is no external certificate infrastructure or
long-term trust store beyond the peer logic maintained by each client.
Third, testing focuses heavily on helper functions and selected protocol
paths rather than large-scale automated interoperability scenarios.

These limitations are acceptable for a course project prototype, but they
also suggest clear next steps for future work.

\section{Future Work}
Several extensions would improve the system:
\begin{itemize}
\item unify feature parity across all three clients,
\item add automated cross-language integration tests,
\item strengthen persistent identity management and trust storage,
\item replace ad hoc command-line interaction with a more robust user
interface, and
\item expand peer discovery and reconnect logic for less controlled
network environments.
\end{itemize}

\section{Conclusion}
SecureDrop demonstrates that a secure file sharing protocol can be
implemented consistently across multiple programming languages while
preserving the core security goals of authentication, confidentiality,
integrity, and origin verification. The project combines Ed25519 identity
keys, X25519 session establishment, AES-GCM encrypted transport, SHA-256
content hashing, and metadata-preserving re-sharing into a coherent
end-to-end design. The resulting system is a strong educational prototype
that highlights both the value and the practical challenges of
cross-language secure systems development.

\end{document}
